\chapter[What the literatures contribute]{What the literatures contribute, and where their assumptions stop}
\label{sim:literature-synthesis}

The simulation problem crosses several literatures because it combines inference from incomplete reports, accounting representation, dynamic corporate finance, sequential information and computational approximation. No inspected source supplies the entire system. The useful synthesis identifies the problem each source actually solves and the additional work needed for AccDSL. A source's notation or software can be valuable without becoming a universal model of a disclosed company.

\section{Inverse problems and incomplete information}
\citet{stuart2010} develops Bayesian inverse problems with explicit forward models, priors and analytical conditions. It is relevant because statements are observations generated by hidden economic histories under a reporting map. The source does not reconstruct company journals or justify a particular prior over contracts. Its contribution is a disciplined distinction between the observation relation and the assumptions needed for inference.

\citet{libkin2014,libkin2016certain} studies incomplete databases and answers across possible worlds. This perspective fits the requirement to preserve alternatives without assigning probabilities. A certain answer holds across admitted worlds; a possible answer has a witness. Translating that logic to finance requires a domain of economic histories and valid accounting relations, not merely null values in a table.

The two approaches are complementary. A feasible family expresses what has not been ruled out. A probability measure expresses relative weight under a model. A Bayesian posterior can be computed on a family when priors and likelihoods are justified, but the family itself should not be discarded when those assumptions are uncertain.

\section{Set identification and prior sensitivity}
\citet[Sections 2.1--2.3]{giacomini2021} separates parameters identified by the observable distribution from remaining directions whose conditional priors are not revised by data. The retained technical reading uses the April 2020 working-paper version associated with the later publication. That version identity matters when citing equations and claims.

For company simulation, the immediate implication is conditional: first determine which hidden features the observation model identifies. A long history of interest expense may leave fixed-versus-floating terms unresolved when rates have not varied. Additional data of the same uninformative kind need not eliminate the ambiguity.

The local history family is not automatically the identified set of the source's statistical model. Finite compatibility with reports is weaker than equality of full observable laws. The manuscript therefore uses report-compatible histories unless the stronger identification argument has been established. This avoids borrowing a technical term to imply more evidence than exists.

\section{Model averaging}
\citet{hoeting1999} surveys Bayesian model averaging, in which uncertainty across models is represented by posterior weights. This can be useful when alternative economic or observation models are explicitly specified. It also makes clear that averaging requires model priors and likelihoods, not branch counts.

The practical difficulty for AccDSL is defining model identity. Ten syntactic encodings of one contract should not receive ten times the weight of a single encoding of another. A prior over economic structures must be invariant to irrelevant representation choices or at least expose its dependence.

Model averaging can also conceal a decision-relevant minority case if the report shows only a mean. Conditional outcomes, tail measures and sensitivity to prior weights remain important. When weights are not defensible, the unweighted family is the appropriate retained representation.

\section{Causal scenarios}
\citet[Sections 3.2--3.4]{pearl2009} distinguishes observational conditioning, interventions and counterfactual reasoning. This distinction matters whenever a financial scenario is described as the effect of a shock. A model fitted to periods with high rates need not identify the effect of changing rates while other mechanisms remain fixed.

AccDSL can support deterministic stress paths without claiming causal identification. Such a run answers a conditional mechanical question. A causal claim additionally needs structural equations, intervention semantics and assumptions about confounding and invariance.

Counterfactual comparisons require a coupling between factual and alternative worlds. Preserving the same hidden history and exogenous innovations is one modeled coupling. Independently drawing a new company for each policy changes the comparison. The causal literature supplies the conceptual distinction; the company-specific structure remains to be justified.

\section{Corporate microsimulation}
\citet{shahnazarian2004} links financial stocks through accounting bridges and separately specifies behavioral relationships for incorporated businesses. The inspected sections include stock movements, investment, disposal, depreciation and tax-related distinctions. The source's Swedish data and tax setting are not modern universal issuer defaults.

Its relevance is methodological: identities constrain movement while behavioral equations determine choices. A simulator can preserve that separation while using different data, accounting rules and policies. Importing estimated coefficients without target-specific evidence would be an unsupported transfer.

The source also illustrates why state includes more than current statement totals. Tax reserves, capital and other accumulated quantities affect future responses. AccDSL's state design should identify such memory by mechanism rather than simply copying a fixed list of variables from one model.

\section{Accounting in system dynamics}
\citet{pierson2020} develops accounting reporting within a system-dynamics model. Accrual, payment and reporting accumulation are separate, and the retained author model exposes equations for cash, receivables, payables and report timing. These are useful implementation concepts.

The source's payables payment rule uses a delay approximation, and its financing and inventory mechanisms have stated limits. A first-order payment delay is not an invoice-level due-date schedule. Reporting values expressed as annualized rates must be converted before comparison with period amounts.

The source therefore contributes a bridge between continuous simulation stocks and discrete reports. It does not establish a complete instrument-level financing engine or empirical validity for a particular company. AccDSL should reuse the separation of clocks while admitting more detailed schedules when the query requires them.

\section{Agent-based stock-flow consistency}
\citet{caiani2016} models heterogeneous agents, loans and capital vintages with an explicit within-period sequence. The retained paper and selected official code show production, credit, payments, taxes and distributions connected through counterpart balance changes. This supports the importance of event ordering and financial counterparts.

The source begins from an artificial economy with its own initialization and calibration. That is not a reconstruction of an observed issuer from disclosures. Its reported simulation properties are author results, not experiments reproduced for AccDSL.

A particularly useful source-to-target warning concerns its operating-cash-flow label: the inspected definition subtracts scheduled principal repayments. That quantity can be appropriate for its credit-demand mechanism, but is not automatically the same as a financial-statement operating cash-flow category. Definitions must travel with values across adapters.

Selected code inspection is also bounded. The retained implementation revision exposes event switches, debt payment calculations and bankruptcy paths. It does not certify every formula, the entire runtime or correspondence with the paper-era executable. Before reuse, zero-rate schedules, timing, numerical domains and accounting adapters need targeted checks.

\section{Financial planning and simultaneous equations}
The retained Gentry working paper studies strategic financial planning from management inputs such as earnings-per-share growth, investment growth, leverage and financing costs \citep{gentry1974wp}. It asks an inverse planning question about returns needed to achieve objectives. Its managerial distributions and independence assumptions are part of that model, not facts about every company.

The retained scan lacks a referenced derivation appendix in the inspected trailing pages, and correspondence with the later journal version has not been established. The manuscript therefore does not rely on an unchecked derivation as a validated algorithm. This is a source-specific limit, not a reason to dismiss the planning approach.

\citet{gentryhoque1990wp} organizes cash-flow components and estimates simultaneous behavioral equations. Its accounting reconciliation, definitions and behavioral relationships perform different jobs. The retained reading identifies a normalization inconsistency in one displayed cash term and records the local algebraic normalization explicitly. It does not silently claim a corrected edition.

The reported sample fits do not establish out-of-sample intervention validity. Firm-level data, a complete instrument matrix and code were not available for local replication. The useful comparison is architectural: a simultaneous system can propose related annual amounts, while a scenario engine still needs dated events, feasibility, information and accounting treatment.

\section{Dynamic corporate finance}
\citet{strebulaev2012} surveys dynamic corporate-finance models with investment, financing frictions and default. The risky-debt section makes debt pricing and equity value interdependent. This is essential for any model claiming endogenous financing terms.

The retained equations also illustrate why source inspection cannot stop at a citation. The face/proceeds convention mismatch in equations (3.26)--(3.27) must be reconciled before implementation. A positive-recovery calculation exposes it; a no-default example can conceal it. The present manuscript uses distinct cash and promise symbols and retains the source qualification.

Other simplifications require equally explicit transfer decisions. Signed net debt can collapse cash and borrowing. Present-valued tax benefits are not actual tax receipts. Aggregate issuance costs do not create legal share rights. A structural model can be internally useful under these assumptions while remaining insufficient for a full accounting simulator.

\section{Bank funding and adjustment}
\citet{halaj2016} models portfolio choices under uncertain funding and constraints. \citet{behn2019} links bank behavior to capital and liquidity restrictions using its own supervisory-data setting and structural assumptions. These sources show that funding is constrained behavior, not an automatic balancing liability.

A corporate issuer does not inherit their bank-specific objective or regulatory measurements. Their estimated parameters are not identified from public company statements merely because the balance-sheet labels resemble those of banks. AccDSL must distinguish accounting, prudential and economic quantities.

The transferable design principle is a policy interface that receives state, information and constraints jointly. Funding availability can change under stress and affect asset decisions. Applying that principle to another sector requires sector-specific rules and evidence, not a relabeled bank equation.

\section{Contract algorithms}
The retained ACTUS technical specification separates schedules, state transitions and payoffs \citep{actus}. This is a promising component boundary for contractual cash flows. It does not supply operating behavior, accounting recognition, consolidation or a justified future risk process.

The inspected specification and service code leave important adapter conventions to be resolved: timezone mapping, midnight interpretation, same-date event order, day-count implementation, currency scaling and risk observations. A service fallback to current wall-clock time plus a horizon is incompatible with reproducible historical simulation unless overridden explicitly.

The selected service requests a particular core dependency but does not establish a fully resolved, matched implementation and dictionary version. No runtime adoption follows from the inspection alone. A conditional algebraic correspondence for a simple borrower loan is useful evidence, but not complete executable conformance.

\section{Sequential decisions under partial observation}
\citet{kaelbling1998} provides the belief-state formulation of planning under partial observation with specified transition and observation probabilities. It explains how information history can be summarized for control. The relevant sections were inspected in the retained author preprint, including rendered pages because the copy is scanned.

AccDSL adds two complications. Some uncertainty is unweighted and cannot be represented as a single probability belief without assumptions. Different actors can have different private information. A public analyst's belief is not necessarily the state on which management acts.

The POMDP formulation is therefore a method for an admitted probabilistic subproblem, not the universal semantics of all company uncertainty. It is especially useful for testing nonanticipativity, signal value and belief updates in finite reference cases.

\section{Exact aggregation and knowledge compilation}
\citet{givan2003} gives equivalence and minimization results for MDPs, including policy-preserving stochastic bisimulation. Its fixed-action counterexample explains why equal outputs under preselected action sequences can be too weak. AccDSL's local extension preserves observations and requested outputs as part of a joint successor law.

\citet{darwiche2002} maps representational tradeoffs among compiled logical languages. It supports sharing and query-aware structure while warning that succinctness and tractable operations trade off. It does not prove that arbitrary financial histories can be compiled cheaply.

Together these sources support a careful answer to state explosion: share common structure, preserve joint constraints and aggregate only under a scoped sufficiency condition. They do not support the claim that replacing a tree with a graph eliminates exponential worst cases.

\section{Filtering and approximation}
\citet{doucet2012tutorial} explains particle filtering, smoothing, degeneracy and conditional analytic integration in the retained tutorial version. Its technical treatment motivates weighted particles and diagnostics when a state-space probability model is specified.

Unknown static contract parameters require additional care. Repeated resampling can remove types without any evidence ruling them out. A high effective sample size after resampling does not guarantee historical diversity. The retained feasible family should therefore remain available alongside the particle approximation.

\citet[Sections 2--4]{dupacova2000preprint} studies probability-metric scenario reduction in the retained 2000 preprint. Theorem 3.1 solves redistribution for a fixed subset; the subsequent subset-selection algorithms are heuristics. The stability results depend on the stochastic program and its regularity, while the reported electrical-load experiment measures distribution distances in its own setting. It does not establish company default accuracy or preserve every adaptive policy. The local derivations in Chapter \ref{sim:approximation} explain the resulting expected-output and policy bounds, and the discontinuous-threshold counterexample shows their limits.

\section{Dynamic ambiguity}
\citet{epstein2002preprint} derives recursive multiple-priors preferences under stated axioms and rectangularity relative to an information filtration. The definition, theorem and relevant proof structure were inspected in the retained September 2002 preprint. The result concerns a particular preference representation.

Its relevance is the danger of using stagewise worst-case recursion with an uncertainty set that restricts cross-time dependence. Rectangular enlargement can change the admitted joint laws. A fixed unknown contract cannot silently become an adversary choosing a new type each period.

AccDSL should offer commitment, original-set and recursive robust formulations only with their semantics stated. The choice can alter policy and should remain visible. Dynamic consistency is not a label that makes every robust calculation equivalent.

\section{Learned economic solvers}
\citet{maliar2021} formulates dynamic economic equations for deep learning, with careful distinctions among residual objectives and sampling constructions. Its definitions and algorithm are relevant to the supplied lecture's proposed solver. The double-sampling derivation in this part shows why the target objective matters before training begins.

\citet{duarte2026} is a newer structural-estimation working paper with jointly solved financing and value components, lagged-network training and local refinement. Its assumptions, state and pricing conventions differ from the lecture model. Its numerical organization is a candidate, not a verified AccDSL implementation or general global-optimality result.

Official code for these learned methods has not been comprehensively audited for this local integration. Consequently, the manuscript makes no code-equivalence, convergence or performance promotion claim. Any implementation decision must add that inspection and a target-specific numerical protocol.

\section{What remains to be researched}
The survey is broad but not exhaustive. Important remaining questions include scenario-reduction guarantees for adaptive company policies, robust filtering with static unidentified parameters, strategic disclosure, multi-creditor restructuring, real-company policy estimation and empirical forecast evaluation. Some require additional primary sources; others require derivations and data specific to the target company.

The earlier Francis--Rowell simultaneous-equation lead remains technically uninspected because attempted routes did not provide the paper. It cannot be declared inferior or irrelevant on that basis. The missing Gentry derivation and source-version questions also remain explicit.

A source gap blocks the unsupported claim, not all independent work. The methods already derived here can be specified and checked while the literature queue continues. This keeps research honest without allowing an inaccessible paper to become a reason to stop every implementation task.

\section{The resulting method proposal}
The synthesis favors a joint history family, dated contract and ownership state, actor-specific information, explicit scenario semantics, feasible policy generation, event-to-accounting translation and query-driven computation. Each component has a distinct literature basis and a distinct local obligation.

No single source validates their composition. The worked cases and proposed tests supply initial evidence for the interfaces, while broader implementation and empirical studies remain necessary. The design should therefore preserve the origin and scope of each assumption instead of presenting the synthesis as an established off-the-shelf theory.

The next chapter turns this method proposal into concrete records and algorithms. Its engineering choices remain replaceable. What must survive replacement is the meaning of evidence, uncertainty, information, financing and result status developed throughout this part.
